\documentclass[../../../include/open-logic-section]{subfiles}

\begin{document}

\olfileid{sth}{ordinals}{ordtype}
\olsection{序数作为序型}

有了替换公理，我们现在在 $\ZFminus$ 中工作，终于可以证明我们一直想要得到的结果：

\begin{thm}\ollabel{thmOrdinalRepresentation}
每个良序都同构于唯一的序数。
\end{thm}

\begin{proof}
设 $\tuple{A, <}$ 是一个良序。根据
\olref[basic]{ordisoidentity}，它至多同构于一个
序数。因此，采用归谬法，假设 $\tuple{A, <}$ 不与\emph{任何}序数同构。
我们首先“使 $\tuple{A, <}$ 尽可能地小”。
具体来说：如果某个真前段 $\tuple{A_a,
<_a}$ 不与任何序数同构，那么就存在一个最小的 $a \in A$
具有此性质；则令 $B = A_a$ 且 $\mathord{\lessdot} =
\mathord{<_a}$。否则，令 $B = A$ 且 $\mathord{\lessdot} =
\mathord{<}$。

根据定义，$B$ 的每个真前段都同构于某个序数，且如前所述该序数唯一。
因此，由替换公理，下述集合存在并且是一个函数：
\[
	f = \Setabs{\tuple{\beta, b}}{b \in B\text{ 且 }\ordeq{\beta}{\tuple{B_b, \lessdot_b}}}
\]
为完成归谬，我们将证明存在某个序数 $\alpha$，使得 $f$ 是从 $\alpha$ 到 $B$ 的同构。

显然 $\ran{f}
= B$。并且根据 \olref[iso]{lemordsegments}，$f$ 保持序关系，
即 $\gamma \in \beta$ 当且仅当 $f(\gamma) \lessdot f(\beta)$。为证明 $\dom{f}$ 是一个序数，根据
\olref[basic]{corordtransitiveord}，只需证
$\dom{f}$ 是传递的即可。取定 $\beta \in \dom{f}$，即
对某个 $b$，$\ordeq{\beta}{\tuple{B_b, \lessdot_b}}$ 成立。若 $\gamma \in
\beta$，则由
\olref[iso]{wellordinitialsegment} 有 $\gamma \in \dom{f}$；
概括得 $\beta \subseteq
\dom{f}$。
\end{proof}

该结果使得我们可以给出以下定义，这也是我们自 \olref[vn]{sec} 以来一直想要给出的：

\begin{defn}
若 $\tuple{A, <}$ 是一个良序，则其序型
$\ordtype{A, <}$ 是使得 $\ordeq{\tuple{A, < }}{\alpha}$ 成立的唯一序数 $\alpha$。
\end{defn}

此外，这个定义还引出两条很好的原理：

\begin{cor}\ollabel{ordtypesworklikeyouwant}
设 $\tuple{A, <}$ 和 $\tuple{B, \lessdot}$ 是良序：
\begin{align*}
	\ordtype{A, <} = \ordtype{B, \lessdot}&\text{ 当且仅当 }\ordeq{\tuple{A, <}}{\tuple{B, \lessdot}}\\
	\ordtype{A, <} \in \ordtype{B, \lessdot}&\text{ 当且仅当 }\ordeq{\tuple{A, <}}{\tuple{B_b, \lessdot_b}}\text{ 对某个 }b \in B
\end{align*}
\end{cor}

\begin{proof}
等式成立依据的是 \olref[basic]{ordisoidentity}。
为证第二个断言，令 $\ordtype{A, <} = \alpha$ 且 $\ordtype{B,
\lessdot} = \beta$，并令 $f \colon \beta \to \tuple {B, \lessdot}$
为同构映射。那么：
\begin{align*}
	\alpha \in \beta&\text{ 当且仅当 }\funrestrictionto{f}{\alpha} \colon \alpha \to B_{f(\alpha)}\text{ 是一个同构}\\
	&\text{ 当且仅当 }\ordeq{\tuple{A, <}}{\tuple{B_{f(\alpha)}, \lessdot_{f(\alpha)}}}\\
	&\text{ 当且仅当 }\ordeq{\tuple{A, <}}{\tuple{B_b, \lessdot_b}}\text{ 对某个 $b \in B$}
\end{align*}
这里分别依据的是 \olref[iso]{ordisounique}、
\olref[iso]{wellordinitialsegment} 和
\olref[basic]{ordissetofsmallerord}。
\end{proof}

\end{document}